\section{Methodology}
\label{sec:methodology}

We now present the Difficulty-Normalized Saturation Index (DNSI), motivated by our key observation: benchmark saturation manifests as entropy compression in improvement patterns. We first define the metric, then justify each design decision.

\subsection{Overview}

DNSI quantifies benchmark saturation by measuring the entropy of performance improvements over time, normalized by task difficulty:
\begin{equation}
\text{DNSI} = \frac{H(\Delta_{\text{observed}})}{H_{\text{expected}}(D)}
\end{equation}
where $H(\Delta_{\text{observed}})$ is the entropy of observed improvement deltas and $H_{\text{expected}}(D)$ is the expected entropy given difficulty proxy $D$.

\textbf{Intuition:} A healthy benchmark exhibits diverse improvements --- many research directions, substantial gains --- producing high improvement entropy. A saturated benchmark shows clustered, incremental improvements --- narrow optimization of similar approaches --- producing low entropy. By normalizing by difficulty, we separate true saturation from benchmark hardness: a 100-class task naturally has different improvement dynamics than a 10-class task.

\subsection{DNSI Computation}

\paragraph{Step 1: Extract SOTA History}
From PapersWithCode or equivalent repositories, we extract the time series of SOTA performance:
\begin{equation}
S = \{(t_1, p_1), (t_2, p_2), \ldots, (t_n, p_n)\}
\end{equation}
where $t_i$ is the submission timestamp and $p_i$ is the reported performance (e.g., top-1 accuracy).

\paragraph{Step 2: Compute Improvement Deltas}
We compute windowed improvement deltas using 6-month aggregation:
\begin{equation}
\Delta_w = \sum_{t_i \in w} (p_i - p_{i-1})^+
\end{equation}
where $w$ indexes 6-month windows and $(x)^+ = \max(0, x)$ ensures we count only improvements.

\textbf{Rationale:} 6-month windowing smooths conference clustering (major venues release results in bursts) while preserving temporal dynamics.

\paragraph{Step 3: Compute Improvement Entropy}
We discretize deltas into bins and compute Shannon entropy:
\begin{equation}
H(\Delta) = -\sum_{b} p_b \log_2 p_b
\end{equation}
where $p_b$ is the proportion of deltas falling in bin $b$.

\paragraph{Step 4: Difficulty Normalization}
We normalize by expected entropy given task difficulty:
\begin{equation}
H_{\text{expected}}(D) = \log_2(D)
\end{equation}
where $D$ is a difficulty proxy. For image classification, we use the number of classes:
\begin{equation}
D_{\text{vision}} = N_{\text{classes}}
\end{equation}

\textbf{Rationale:} A 100-class benchmark allows more diverse improvement patterns than a 10-class benchmark. Without normalization, raw entropy conflates saturation with task complexity.

\paragraph{Final DNSI Formula}
\begin{equation}
\text{DNSI} = \frac{H(\Delta_{\text{observed}})}{\log_2(N_{\text{classes}})}
\end{equation}

\textbf{Interpretation:}
\begin{itemize}
    \item DNSI $\approx 1.0$: Improvement entropy matches expected diversity --- healthy benchmark
    \item DNSI $< 0.5$: Improvement entropy significantly below expected --- saturated benchmark
    \item DNSI $> 1.0$: More diversity than expected --- rapidly evolving benchmark
\end{itemize}

\subsection{Design Decisions}

\textbf{Why 6-Month Windows?} Major ML conferences (NeurIPS, ICML, ICLR, ACL, CVPR) cluster submissions, creating artificial periodicity in raw SOTA histories. 6-month windows smooth this clustering while preserving the overall saturation signal.

\textbf{Why Class Count as Difficulty Proxy?} For classification tasks, the number of classes provides an information-theoretic bound on output complexity. Alternative proxies (dataset size, human baseline, theoretical bounds) may improve normalization for non-classification tasks; we leave this extension to future work.

\textbf{Minimum History Requirements:} We require $>15$ SOTA entries over $>3$ years to ensure sufficient data for reliable entropy estimation.
